\chapter{成稿体裁与两份文件}
\label{ch:deliver}

本章登记正文各块装的是哪一类对象，以及同一正文如何分成两份文件。页数是这些对象写开之后的结果。重复段落、未运行套件的示意曲线和把同一中位数换一种说法再写一遍，都不增加技术内容。

\section{250 页所装的内容}

\paragraph{要判定的量}
每一块篇幅对应的是哪一组已经定义的对象：方程、协议、证明、图位或日志。对象没写开时，那一块就短；对象写开之后，不把空图位算成结果。

\paragraph{继承的文献层与不继承的层}
论证顺序继承技术专著里“先定义、再方程、再可否定陈述、最后才是数值”的写法。不继承把篇幅本身当成质量指标。不继承把大纲里的页数预算抄进摘要冒充已经完成的测量。

下表是各块所装对象的清单。右列是规划时为这些对象留出的篇幅，不是对排版软件输出页数的事后声明。实际页数以编译日志为准。

\begin{table}[ht]
\centering
\caption{正文各块所装的对象}
\begin{tabularx}{\textwidth}{lXr}
\toprule
部分 & 写开的对象 & 规划页 \\
\midrule
问题与闭环 & 信息赤字、五条性质、四个判断、六步完整性 & 16 \\
已有能力 & 物理描述、装置实际闭合的段、不构成新增环节的做法 & 与问题并列 \\
文献分层 & 每条来源继承的方程结构与明确不继承的边界 & 32 \\
架构 & 五层、四个时间尺度、模块合同、腔体最小执行 & 18 \\
多尺度模型 & 全场、能量与三相、界面、反应进度、决策流形、两节点、闭合、无量纲 & 40 \\
估计 & 观测算子、病态、灵敏度、Fisher、跟踪器与集合规格的差别 & 28 \\
对偶与探针 & 信息目标、探针库、损伤上界、在线特例、对偶目标 & 26 \\
因果 & 结构方程、干预与相关、路径 A/B/C、混杂、未运行的黑箱对照 & 22 \\
演化 & 算子分工、贝叶斯更新、一次烹饪内的一次方差更新、回退 & 20 \\
控制与安全 & 品质泛函、机会约束的规格、屏障、降级、加热器这一路硬测量 & 24 \\
试验协议 & 双模型、四套套件、主因素、指标、交叉表 & 18 \\
命题证明 & 六条命题的证明与断言常数 & 16 \\
图件 & 88 个图位的坐标、来源和不能支持的句子 & 42 \\
结果 & 只含已落盘的 D 图与交叉表 & 28 \\
实施例 & 三组步骤、同一腔体上尚未同时安装的量 & 16 \\
\midrule
合计 & 规划对象，不是排版承诺 & 约 342 \\
\bottomrule
\end{tabularx}
\end{table}

结果块若某一坐标没有日志，正文保留坐标和图位，不把空白图位计入“已经得到的结果”。证明、文献边界和离散格式不因为结果块变薄就删掉。删掉之后，剩下的中位数没有适用条件。

这一份正文在第二次编译里写出 170 页。规划合计约 342。差额来自尚未运行的对象，以及写开之后仍然短于规划的层。不把差额写成待补的空白页。

本次已经落盘、因而可以进入结果块的对象是：S1 的 4860 条组合乘四种策略、五张由该日志画出的图、六条命题在其断言常数下的通过记录、一条种子 7 的轨迹、以及 140 条清洁 GAB 拟合。尚未落盘、因而不进入结果句的对象是：S0 的网格减半、S2 的遮挡与噪声扫描、S3 的算子训练、T2 的二维脂肪几何、T3 的边界库、集合滤波、CasADi 求解和 FNO 与 DeepONet 的误差对比。

\paragraph{本节不能推出的结论}
规划页数大于编译页数，只说明有的对象还没写到该篇幅，不说明缺页的部分可以用重复段落补上。编译页数大于某一章的规划，只说明公式和长表占了纵向空间，不说明多出来的页是新的实验。

\section{论证顺序}

\paragraph{要判定的量}
一章内部的四段是否齐全，以及摘要里的定量句是否指向一幅已经画出的数据图。

每一章内部固定为四段。

\begin{enumerate}
\item 本章或本节要判定的量。判定失败时读者会把哪个错误句子当成结论，要在这里写明。
\item 继承的文献层与不继承的层。继承指后文直接使用该来源的方程结构。不继承指该来源的具体边界、物料数或装置条款不能被写成本文的测定。
\item 方程、算法或协议。已运行的离散格式与尚未运行的规格必须分成两句。
\item 该节不能推出的结论。这一段缩小句子的主语，不重复上一段的推导。
\end{enumerate}

摘要和每章末尾的定量句遵守下面的指向规则。

\begin{itemize}
\item 有 D 图或交叉表的数，句子里写出图号或表号。本次可以指向的是 D1、D6、D7、D10、D25 和日志交叉表。
\item 没有 D 图的陈述放在定义或命题里，并使用条件句：若假设成立，则结论成立。
\item 设计线 \(0.85\) 只作为 \(e_{\alpha}\) 的零点。它不是认证阈值。
\item 19440 行是同一数值核在因子格子上的样本。句子的谓语用“样本”，不用“次实验”。
\item 浏览器副本的中位数不进入主表。主表只读战役日志。
\end{itemize}

\paragraph{本节不能推出的结论}
四段写全，不能代替命题的证明，也不能代替日志。顺序只防止把效果句写在定义之前。

\section{同一正文的两份文件}

\paragraph{要判定的量}
PDF 与 Word 是否同一套章、节、公式编号和图号。两份文件若各改各的公式，后文的交叉引用就不再指同一个对象。

\begin{itemize}
\item PDF 由本书的 TeX 源生成，公式编号以 PDF 为准。
\item Word 由同一正文导出，用于批注。图号、表号和公式编号与 PDF 相同。
\item 转换若改变断行或公式外观，不另起一套编号。
\item 两份文件使用同一标题、同一参考文献表。正文不印版本号。
\item 演示字段以第 \ref{ch:demo} 章为准，不在两份文件里各写一套列名。
\end{itemize}

技术修改先进入 TeX 源，再同步到 Word。Word 不是第二套技术内容的编辑处。学生身份、班级和提交截止时刻不属于技术对象，不写入正文。

\paragraph{本节不能推出的结论}
两份文件编号一致，不能推出其中的策略排序已经在异构真值上成立。编号一致只约束引用，不提高模型的保真度。

\section{演示所要求的计算形态}

\paragraph{要判定的量}
页面上的温度和路径条，是从数值核的状态流读出来的，还是在着色器里另算了一套差分。

规格把计算和界面分开。

\begin{itemize}
\item 真值积分、估计更新、探针时序和加热削减放在可测试的数值核中。输入输出字段与第 \ref{ch:demo} 章一致。
\item 数值核有独立于界面的断言。命题 \ref{prop:steady}、\ref{prop:pulse}、\ref{prop:budget}、\ref{prop:shield} 各有一条可重复核对。脉冲断言使用的是“两段滞后大于单段滞后加 1”的绝对准则。相对 1.15 倍的准则在同一组响应上不成立，不作为通过条件。
\item 界面读取数值核的状态流。界面不另写一套导热差分。
\item 下载文件由数值核序列化。浏览器里为了观看而保存的副本，若列集合或中位数与战役日志不同，不进入主表。
\item 多孔介质离线结果若以后接入，列名保持不变。判据不改。
\end{itemize}

已运行的核是十六条控制体上的工程通量，加上标量跟踪器。它满足“可重复”和“断言独立于界面”。它不满足“积分器是多孔介质真值”和“估计器不读取真值温差”。种子 7 的演示初温 \(8^\circ\mathrm{C}\) 不在 S1 的初温格子 \(\{5,20\}^\circ\mathrm{C}\) 里，因此那一条轨迹用于核对更新规则，不进入 4860 的分母。

\paragraph{本节不能推出的结论}
断言通过，不能写成 S1 的中位数已经证明了该命题的定量常数。断言常数与战役窗口是两套数：冷点上界断言用 \(\tau_2=30\,\mathrm{s}\)、\(\tau=3\,\mathrm{s}\)，战役探针窗口是 40 s。
